\chapter{Hilbert Spaces}

\begin{center}\itshape
Pythagoras's theorem holds for functions. A function is a point, its Fourier coefficients are its coordinates, and the best
approximation of a point by a subspace is the foot of a perpendicular. How far can the geometry of the plane be pushed?
\end{center}

Around 1900, the hot topic in G\"ottingen was integral equations. Ivar Fredholm, a Swede who also worked in insurance and later as an
actuary, had shown in 1900 how to solve $f(x)-\int_0^1K(x,y)f(y)\,dy=g(x)$ by treating it as the limit of a linear system
with infinitely many unknowns. David Hilbert heard about it in the winter of 1900--1901 from his student Holmgren, dropped what he was
doing, and spent the next decade on it. In six long papers (1904--1910) he invented the spectral theory of ``quadratic forms in
infinitely many variables'', the variables being the coefficients $x_1,x_2,\dots$ of a function in an orthonormal system, with
$\sum x_n^2<\infty$.

It was Hilbert's student Erhard Schmidt who, in 1908, saw that all this was \emph{geometry}. In Schmidt's paper the sequences with
$\sum|x_n|^2<\infty$ have a length, two of them can be orthogonal, Pythagoras's theorem holds, and a point has a projection on a
closed subspace. The year before, in 1907, Frigyes Riesz and Ernst Fischer had proved independently that $L^2([a,b])$ is complete, so that
\emph{every} square-summable sequence is the sequence of Fourier coefficients of some function. Functions and sequences were the same
space, seen in two coordinate systems. When the young John von Neumann axiomatised the space in 1929 for quantum mechanics, he called it
\emph{Hilbert space}. A well-known story, which we could not trace to a first-hand source and which may be legend, has Hilbert attending a lecture
on quantum mechanics a little later and asking the speaker what, after all, a Hilbert space is.

This chapter answers him. It proves that the closest point of a closed convex set always exists, that the dual of a Hilbert space is
itself, and that every separable Hilbert space is a copy of $\ell^2$. Fourier series find their natural home, von Neumann's ergodic
theorem shows that averages along an orbit converge, his proof of the Radon--Nikodym theorem takes ten lines, and spaces of functions
with a kernel turn into the most practical tool of machine learning.

\section{Inner products}

An \emph{inner product} on a vector space $H$ over $\mathbb K$ is a map $\langle\cdot,\cdot\rangle:H\times H\to\mathbb K$, linear in the first variable, with $\langle y,x\rangle=\overline{\langle x,y\rangle}$ and $\langle x,x\rangle>0$ for
$x\ne0$. It is conjugate-linear in the second variable. Put $\|x\|=\langle x,x\rangle^{1/2}$. Expanding $\|x+ty\|^2\ge0$ gives the \emph{Cauchy--Schwarz inequality} $|\langle x,y\rangle|\le\|x\|\|y\|$ (Volume~0), and with it
$\|x+y\|^2=\|x\|^2+2\operatorname{Re}\langle x,y\rangle+\|y\|^2\le(\|x\|+\|y\|)^2$: $\|\cdot\|$ is a norm. A \emph{Hilbert space} is an inner product space complete for this norm.

The examples are $\mathbb K^n$, $\ell^2$ with $\langle x,y\rangle=\sum_nx_n\overline{y_n}$, and above all $L^2(\mu)$ with $\langle f,g\rangle=\int f\overline g\,d\mu$, complete by the Riesz--Fischer theorem (Volume~6, Theorem~9.4). Two vectors are
\emph{orthogonal}, $x\perp y$, if $\langle x,y\rangle=0$; then $\|x+y\|^2=\|x\|^2+\|y\|^2$. Adding the expansions of $\|x+y\|^2$ and $\|x-y\|^2$ gives the \emph{parallelogram law}
\[\|x+y\|^2+\|x-y\|^2=2\|x\|^2+2\|y\|^2,\]
which says that the diagonals of a parallelogram carry as much squared length as its four sides. It characterises the norms that come from an inner product (Exercise~\ref{ex:jvn}), and it is the only property of Hilbert
spaces that the next theorem uses.

\section{The projection theorem}

\begin{theorem}[projection on a closed convex set]\label{thm:projection}
Let $C$ be a non-empty closed convex subset of a Hilbert space $H$ and $x\in H$. There is a unique $p=P_C(x)\in C$ with $\|x-p\|=d(x,C)$. It is characterised by
\[p\in C\quad\text{and}\quad\operatorname{Re}\langle x-p,c-p\rangle\le0\quad\text{for all }c\in C,\]
and $\|P_C(x)-P_C(y)\|\le\|x-y\|$.
\end{theorem}
\begin{proof}
\emph{Existence.} Let $d=d(x,C)$ and $c_n\in C$ with $\|x-c_n\|^2\le d^2+\frac1n$. The parallelogram law applied to $x-c_n$ and $x-c_m$ gives
\[\|c_n-c_m\|^2=2\|x-c_n\|^2+2\|x-c_m\|^2-4\Bigl\|x-\frac{c_n+c_m}2\Bigr\|^2\le2\Bigl(d^2+\frac1n\Bigr)+2\Bigl(d^2+\frac1m\Bigr)-4d^2,\]
because $\frac{c_n+c_m}2\in C$ by convexity. So $(c_n)$ is Cauchy; its limit $p$ is in $C$ (closed) and $\|x-p\|=d$. The same computation with two minimisers $p,p'$ in place of $c_n,c_m$ gives $\|p-p'\|^2\le0$: uniqueness.

\emph{Characterisation.} If $p$ is the minimiser and $c\in C$, then $p+t(c-p)\in C$ for $t\in]0,1]$, so
$\|x-p\|^2\le\|x-p-t(c-p)\|^2=\|x-p\|^2-2t\operatorname{Re}\langle x-p,c-p\rangle+t^2\|c-p\|^2$; divide by $t$ and let $t\to0$. Conversely, if $p\in C$ satisfies the inequality, then for $c\in C$,
$\|x-c\|^2=\|x-p\|^2+\|p-c\|^2-2\operatorname{Re}\langle x-p,c-p\rangle\ge\|x-p\|^2$.

\emph{Contraction.} With $p=P_C(x)$, $q=P_C(y)$, the characterisation gives $\operatorname{Re}\langle x-p,q-p\rangle\le0$ and $\operatorname{Re}\langle y-q,p-q\rangle\le0$. Adding, $\operatorname{Re}\langle(x-y)-(p-q),q-p\rangle\le0$, that is
$\|p-q\|^2\le\operatorname{Re}\langle x-y,p-q\rangle\le\|x-y\|\|p-q\|$.
\end{proof}

\begin{figure}[ht]
\centering
\fig{ch02-projection}
\caption{The projection $p$ of $x$ on a closed convex set $C$: the angle between $x-p$ and $c-p$ is obtuse for every $c\in C$. On a closed subspace the angle is right, and $p$ is the foot of the perpendicular.}
\label{fig:projection}
\end{figure}

When $C=M$ is a closed subspace, $c-p$ runs over all of $M$, and the inequality applied to $\pm m$ and $\pm im$ becomes $\langle x-p,m\rangle=0$: $P_M(x)$ is the unique $p\in M$ with $x-p\perp M$. The \emph{orthogonal}
$M^\perp=\{y:y\perp M\}$ is a closed subspace (an intersection of kernels of continuous forms).

\begin{corollary}\label{cor:orthdecomp}
If $M$ is a closed subspace, then $H=M\oplus M^\perp$, and $P_M$ is linear with $P_M^2=P_M$, $\|P_M\|\le1$ and $\langle P_Mx,y\rangle=\langle x,P_My\rangle$. For any subspace $N$, $(N^\perp)^\perp=\overline N$; so $N$ is dense if and only if
$N^\perp=\{0\}$.
\end{corollary}
\begin{proof}
$x=P_Mx+(x-P_Mx)$ with $x-P_Mx\in M^\perp$, and $M\cap M^\perp=\{0\}$ because $m\perp m$ forces $m=0$. Linearity follows from the characterisation (if $x-p\perp M$ and $y-q\perp M$, then $(\lambda x+y)-(\lambda p+q)\perp M$), and
$\langle P_Mx,y\rangle=\langle P_Mx,P_My\rangle=\langle x,P_My\rangle$. For any $N$, $N^\perp=\overline N^\perp$ by continuity of the inner product; clearly $\overline N\subseteq(\overline N^\perp)^\perp$, and if $x\perp\overline N^\perp$, then $x-P_{\overline N}x$ is in
$\overline N^\perp$ and orthogonal to itself, so $x\in\overline N$.
\end{proof}

\begin{think}
In $\R^2$ with the norm $\|(a,b)\|_\infty=\max(|a|,|b|)$, how many points of the line $b=0$ are closest to $(0,1)$? Which step of the proof of Theorem~\ref{thm:projection} fails? (Exercise~\ref{ex:noclosest} shows that a closest
point may not exist at all in a Banach space.)
\end{think}

\section{The Riesz representation theorem}

Every $y\in H$ defines a bounded linear form $x\mapsto\langle x,y\rangle$ of norm $\|y\|$ (Cauchy--Schwarz, with equality at $x=y$). There are no others. This was proved by Fr\'echet and by Riesz in 1907, in the same issue of the
\emph{Comptes rendus}.

\begin{theorem}[Riesz--Fr\'echet]\label{thm:rieszrep}
For every $\varphi\in H^*$ there is a unique $y\in H$ with $\varphi(x)=\langle x,y\rangle$ for all $x$, and $\|\varphi\|=\|y\|$. The map $\varphi\mapsto y$ is a conjugate-linear isometric bijection $H^*\to H$.
\end{theorem}
\begin{proof}
If $\varphi=0$, take $y=0$. Otherwise $M=\ker\varphi$ is a closed subspace different from $H$, so $M^\perp\ne\{0\}$ (Corollary~\ref{cor:orthdecomp}); take $z\in M^\perp$ with $\|z\|=1$. For every $x$, $\varphi(x)z-\varphi(z)x\in M$, so it is orthogonal to $z$:
$\varphi(x)-\varphi(z)\langle x,z\rangle=0$. Hence $\varphi(x)=\langle x,y\rangle$ with $y=\overline{\varphi(z)}z$. Uniqueness: if $\langle x,y-y'\rangle=0$ for all $x$, take $x=y-y'$.
\end{proof}

In particular Hilbert spaces are reflexive, and their bounded sets are weakly sequentially compact.

\begin{corollary}\label{cor:weakcompact}
Every bounded sequence in a Hilbert space has a weakly convergent subsequence. If $x_n\rightharpoonup x$ and $\|x_n\|\to\|x\|$, then $x_n\to x$ in norm.
\end{corollary}
\begin{proof}
Let $\|x_n\|\le M$ and $H_0$ the closed span of the $x_n$, a separable Hilbert space. By diagonal extraction, a subsequence $(y_k)$ has $\langle y_k,x_n\rangle$ convergent for every $n$, hence $\langle y_k,z\rangle$ convergent for every $z$ in
the span of the $x_n$, and then for every $z\in H_0$ (approximate $z$ within $\varepsilon$ and use $|\langle y_k,z-z'\rangle|\le M\varepsilon$). The limit $\overline{\varphi(z)}$ defines $\varphi\in H_0^*$ with $\|\varphi\|\le M$; by
Theorem~\ref{thm:rieszrep} in $H_0$ there is $x\in H_0$ with $\langle y_k,z\rangle\to\langle x,z\rangle$ for $z\in H_0$. For $z\in H$, $\langle y_k,z\rangle=\langle y_k,P_{H_0}z\rangle\to\langle x,P_{H_0}z\rangle=\langle x,z\rangle$. So $y_k\rightharpoonup x$.
For the second statement, $\|x_n-x\|^2=\|x_n\|^2-2\operatorname{Re}\langle x_n,x\rangle+\|x\|^2\to\|x\|^2-2\|x\|^2+\|x\|^2=0$.
\end{proof}

\section{Orthonormal bases}

A family $(e_i)_{i\in I}$ is \emph{orthonormal} if $\langle e_i,e_j\rangle$ is $1$ for $i=j$ and $0$ otherwise. For a finite orthonormal family $e_1,\dots,e_n$, the projection on its span is $\sum_k\langle x,e_k\rangle e_k$ (the difference with $x$ is
orthogonal to every $e_k$), and Pythagoras gives $\|x\|^2=\sum_{k\le n}|\langle x,e_k\rangle|^2+\|x-\sum_{k\le n}\langle x,e_k\rangle e_k\|^2$. Hence \emph{Bessel's inequality}: $\sum_i|\langle x,e_i\rangle|^2\le\|x\|^2$ for every orthonormal
family, finite or not. In particular only countably many $\langle x,e_i\rangle$ are non-zero. We state the next theorem for sequences; the case of uncountable families is the same with sums over arbitrary index sets.

\begin{theorem}[orthonormal bases]\label{thm:onb}
Let $(e_n)_{n\ge1}$ be an orthonormal sequence in a Hilbert space $H$. For every $x$, the series $\sum_n\langle x,e_n\rangle e_n$ converges, to the projection of $x$ on the closed span of the $e_n$. The following are equivalent:
\begin{enumerate}
\item[(a)] the span of the $e_n$ is dense in $H$;
\item[(b)] if $\langle x,e_n\rangle=0$ for all $n$, then $x=0$;
\item[(c)] $x=\sum_n\langle x,e_n\rangle e_n$ for every $x\in H$;
\item[(d)] (Parseval) $\|x\|^2=\sum_n|\langle x,e_n\rangle|^2$ for every $x\in H$.
\end{enumerate}
Such a sequence is a \emph{Hilbert basis}. Moreover, for every $(c_n)\in\ell^2$ the series $\sum_nc_ne_n$ converges, so that $x\mapsto(\langle x,e_n\rangle)_n$ is an isometric bijection from $H$ onto $\ell^2$ when the $e_n$ form a Hilbert basis.
\end{theorem}
\begin{proof}
If $\sum_n|c_n|^2<\infty$, the partial sums $s_N=\sum_{n\le N}c_ne_n$ satisfy $\|s_N-s_M\|^2=\sum_{M<n\le N}|c_n|^2$, so they form a Cauchy sequence and converge. By Bessel, this applies to $c_n=\langle x,e_n\rangle$; the sum $s$ satisfies
$\langle x-s,e_n\rangle=0$ for all $n$ (continuity of the inner product), so $x-s$ is orthogonal to the closed span $M$ of the $e_n$, and $s\in M$: $s=P_Mx$. (a)$\Rightarrow$(b): such an $x$ is orthogonal to a dense subspace, hence to
itself. (b)$\Rightarrow$(c): $x-s$ is orthogonal to all $e_n$. (c)$\Rightarrow$(d): $\|s_N\|^2=\sum_{n\le N}|\langle x,e_n\rangle|^2$ and $\|s_N\|\to\|x\|$. (d)$\Rightarrow$(a): $\|x-s_N\|^2=\|x\|^2-\sum_{n\le N}|\langle x,e_n\rangle|^2\to0$, and $s_N$
is in the span.
\end{proof}

\begin{corollary}\label{cor:separable}
Every infinite-dimensional separable Hilbert space has a Hilbert basis, and is therefore isometrically isomorphic to $\ell^2$.
\end{corollary}
\begin{proof}
Let $(x_k)$ be dense. Discard every $x_k$ that is in the span of the previous ones, and apply the Gram--Schmidt process to the rest: $f_1=x_{k_1}$, $f_n=x_{k_n}-\sum_{j<n}\langle x_{k_n},e_j\rangle e_j$, $e_n=f_n/\|f_n\|$. The $e_n$ are orthonormal
and have the same span as the $x_k$, which is dense; Theorem~\ref{thm:onb} concludes.
\end{proof}

Without separability, Zorn's lemma gives a maximal orthonormal family, which satisfies (b), and $H$ is isometric to $\ell^2(I)$ for some set $I$. All Hilbert spaces of the same ``dimension'' are the same. This is
the answer to Hilbert's question: up to isometry, the Hilbert space is $\ell^2$, and the content of a theorem about a particular $L^2$ is the choice of a good basis.

\section{Fourier series and orthogonal polynomials}

On $L^2(\mathbb T)=L^2([-\pi,\pi])$ with the normalised inner product $\langle f,g\rangle=\frac1{2\pi}\int_{-\pi}^\pi f\overline g$, the functions $e_n(t)=e^{int}$, $n\in\Z$, are orthonormal, and $\langle f,e_n\rangle=\hat f(n)$ is the $n$-th Fourier
coefficient (Volume~4, Chapter~6).

\begin{theorem}\label{thm:fourierl2}
The trigonometric system $(e^{int})_{n\in\Z}$ is a Hilbert basis of $L^2(\mathbb T)$. Hence, for every $f\in L^2(\mathbb T)$, $\sum_{|n|\le N}\hat f(n)e^{int}\to f$ in $L^2$, $\frac1{2\pi}\int|f|^2=\sum_n|\hat f(n)|^2$, and every
$(c_n)\in\ell^2(\Z)$ is the sequence of Fourier coefficients of a unique $f\in L^2(\mathbb T)$.
\end{theorem}
\begin{proof}
By Theorem~\ref{thm:onb} (for an orthonormal family indexed by $\Z$) it suffices that trigonometric polynomials are dense. Let $f\in L^2$ and $\varepsilon>0$. By Volume~6, Theorem~9.7, there is a continuous $g$ with compact support in
$]-\pi,\pi[$ and $\|f-g\|_2<\varepsilon$; it extends to a continuous $2\pi$-periodic function. By the trigonometric Weierstrass theorem (Volume~6, Corollary~4.20(b)) there is a trigonometric polynomial $P$ with
$\|g-P\|_\infty<\varepsilon$, and $\|g-P\|_2\le\|g-P\|_\infty$ for the normalised norm. So $\|f-P\|_2<2\varepsilon$.
\end{proof}

Volume~4 proved Parseval's identity for piecewise continuous functions; the new content is that it holds for every square-integrable function, and the converse statement, which needs completeness: any square-summable sequence of
coefficients, however wild, is the spectrum of an honest function. A sequence like $c_n=\frac1{|n|^{0.51}}$ is the Fourier series of a function in $L^2$ whose Fourier series converges absolutely at no point.

The same method works on an interval with polynomials. On $L^2([-1,1])$, Gram--Schmidt applied to $1,x,x^2,\dots$ produces orthogonal polynomials, which turn out to be the \emph{Legendre polynomials}
\[P_n(x)=\frac1{2^nn!}\frac{d^n}{dx^n}(x^2-1)^n,\qquad P_0=1,\ P_1=x,\ P_2=\tfrac12(3x^2-1),\ P_3=\tfrac12(5x^3-3x).\]

\begin{proposition}\label{prop:legendre}
The $P_n$ are orthogonal in $L^2([-1,1])$ with $\int_{-1}^1P_n^2=\frac2{2n+1}$, and the $\sqrt{n+\frac12}\,P_n$ form a Hilbert basis.
\end{proposition}
\begin{proof}
Let $u=(x^2-1)^n$, of degree $2n$; it vanishes to order $n$ at $\pm1$, so $u^{(k)}(\pm1)=0$ for $k<n$. If $q$ is a polynomial, $k$ integrations by parts move all derivatives off $u$ with no boundary terms:
$\int_{-1}^1u^{(n)}q=(-1)^n\int_{-1}^1u\,q^{(n)}$. If $\deg q<n$, this is $0$: $P_n$ is orthogonal to all polynomials of lower degree, hence to $P_m$ for $m<n$. With $q=u^{(n)}$, $q^{(n)}=u^{(2n)}=(2n)!$, so
\[\int_{-1}^1(u^{(n)})^2=(2n)!\int_{-1}^1(1-x^2)^n\,dx=(2n)!\,\frac{2^{2n+1}(n!)^2}{(2n+1)!},\]
by the Beta integral (substitute $x=2t-1$; Volume~4, Chapter~5), and dividing by $(2^nn!)^2$ gives $\frac2{2n+1}$. Polynomials are dense in $C([-1,1])$ for the sup norm (Weierstrass, Volume~6, Corollary~4.20(a)), hence in
$L^2([-1,1])$ as in Theorem~\ref{thm:fourierl2}, and $P_0,\dots,P_n$ span the polynomials of degree at most $n$.
\end{proof}

\begin{figure}[ht]
\centering
\fig{ch02-legendre}
\caption{The Legendre polynomials $P_1,\dots,P_5$ on $[-1,1]$. Each $P_n$ has $n$ simple zeros in $]-1,1[$ and is orthogonal to all polynomials of lower degree. The zeros are the nodes of Gauss's quadrature.}
\label{fig:legendre}
\end{figure}

Other weights give other families: Hermite polynomials for $e^{-x^2}$ on $\R$ (whose completeness needs the Fourier transform, Chapter~18), Laguerre for $e^{-x}$ on $[0,\infty[$, Chebyshev for $(1-x^2)^{-1/2}$. Each is
the right basis for some problem, and the projection theorem says what they all do: $\sum_{k\le n}\langle f,e_k\rangle e_k$ is the \emph{best} approximation of $f$ by a polynomial of degree $n$ in the mean-square sense.

\section{Averages along orbits}

Volume~4 proved Weyl's theorem: if $\alpha$ is irrational, the points $n\alpha$ modulo $1$ are equidistributed. Here is the version in several dimensions and, behind it, a theorem of von Neumann that belongs to Hilbert space
geometry. Write $\mathbb T^d=\R^d/\Z^d$, with Lebesgue measure on $[0,1[^d$, and $e_k(x)=e^{2\pi ik\cdot x}$ for $k\in\Z^d$. By the argument of Theorem~\ref{thm:fourierl2}, with Stone--Weierstrass in $d$ variables (Volume~6,
Theorem~4.19(b): the $e_k$ separate the points of $\mathbb T^d$), trigonometric polynomials are dense in $C(\mathbb T^d)$, and the $e_k$ form a Hilbert basis of $L^2(\mathbb T^d)$.

\begin{theorem}[Kronecker--Weyl]\label{thm:kroneckerweyl}
Let $\alpha\in\R^d$ be such that $1,\alpha_1,\dots,\alpha_d$ are linearly independent over $\Q$. For every continuous $f$ on $\mathbb T^d$, uniformly in $x$,
\[\frac1N\sum_{n=0}^{N-1}f(x+n\alpha)\longrightarrow\int_{\mathbb T^d}f.\]
In particular every orbit $\{x+n\alpha\}$ is dense in $\mathbb T^d$, and equidistributed: the proportion of $n<N$ with $x+n\alpha\in B$ tends to the volume of $B$, for every box $B$.
\end{theorem}
\begin{proof}
For $f=e_k$ with $k\ne0$, $k\cdot\alpha\notin\Z$ by the independence hypothesis, so $e^{2\pi ik\cdot\alpha}\ne1$ and the geometric sum gives
\[\Bigl|\frac1N\sum_{n<N}e_k(x+n\alpha)\Bigr|=\frac1N\Bigl|\frac{e^{2\pi iNk\cdot\alpha}-1}{e^{2\pi ik\cdot\alpha}-1}\Bigr|\le\frac2{N|e^{2\pi ik\cdot\alpha}-1|}\to0=\int e_k,\]
uniformly in $x$. For $k=0$ both sides are $1$. By linearity the theorem holds for trigonometric polynomials. Given $f$ continuous and $\varepsilon>0$, take $P$ with $\|f-P\|_\infty<\varepsilon$; then averages and integrals of $f$
and $P$ differ by at most $\varepsilon$, so $\limsup_N\sup_x|\frac1N\sum_{n<N}f(x+n\alpha)-\int f|\le2\varepsilon$. For a box $B$ and $\varepsilon>0$, there are continuous $g\le\mathbf 1_B\le h$ with $\int(h-g)<\varepsilon$; the proportion of
visits lies between the averages of $g$ and of $h$, whose limits are within $\varepsilon$ of $\lambda(B)$. Density follows, since every open set contains a small box.
\end{proof}

\begin{figure}[ht]
\centering
\fig{ch02-kronecker}
\caption{The first $600$ points of the orbit $n(\sqrt2,\sqrt3)$ modulo $1$ in the unit square. Since $1,\sqrt2,\sqrt3$ are independent over $\Q$, the orbit fills the torus evenly (Theorem~\ref{thm:kroneckerweyl}).}
\label{fig:kronecker}
\end{figure}

The proof used the specific structure of the torus. Von Neumann found in 1931 that the convergence of averages holds in $L^2$ for \emph{any} measure-preserving transformation, and that it is a statement about operators of norm
at most $1$ on a Hilbert space.

\begin{theorem}[von Neumann's mean ergodic theorem]\label{thm:meanergodic}
Let $T$ be a bounded operator on a Hilbert space with $\|T\|\le1$, and $P$ the orthogonal projection on $\ker(I-T)=\{x:Tx=x\}$. For every $x$,
\[A_Nx=\frac1N\sum_{n=0}^{N-1}T^nx\longrightarrow Px\quad\text{in norm}.\]
\end{theorem}
\begin{proof}
Recall the \emph{adjoint} $T^*$, defined by $\langle Tx,y\rangle=\langle x,T^*y\rangle$ (by Theorem~\ref{thm:rieszrep}, $y\mapsto T^*y$ is well defined, linear, and $\|T^*\|=\|T\|$). If $Tx=x$, then
\[\begin{aligned}\|T^*x-x\|^2&=\|T^*x\|^2-2\operatorname{Re}\langle T^*x,x\rangle+\|x\|^2\\&=\|T^*x\|^2-2\operatorname{Re}\langle x,Tx\rangle+\|x\|^2\eqbr\le\|x\|^2-2\|x\|^2+\|x\|^2=0,\end{aligned}\]
so $T^*x=x$; by symmetry ($T^{**}=T$) the fixed points of $T$ and $T^*$ coincide. Now $y\perp(I-T)z$ for all $z$ means $(I-T^*)y=0$, so the orthogonal of the range $R$ of $I-T$ is $\ker(I-T^*)=\ker(I-T)$, and by
Corollary~\ref{cor:orthdecomp}, $H=\ker(I-T)\oplus\overline R$. On $\ker(I-T)$, $A_Nx=x=Px$. For $y=(I-T)z\in R$, the sum telescopes: $A_Ny=\frac1N(z-T^Nz)$, of norm at most $\frac2N\|z\|$. For $y\in\overline R$ and
$\varepsilon>0$, pick $y'\in R$ with $\|y-y'\|<\varepsilon$; since $\|A_N\|\le1$, $\limsup\|A_Ny\|\le\varepsilon$. So $A_N\to0=P$ on $\overline R$.
\end{proof}

If $\tau:X\to X$ preserves a probability measure $\mu$, then $Tf=f\circ\tau$ is an isometry of $L^2(\mu)$, and the theorem says that the time averages $\frac1N\sum_{n<N}f\circ\tau^n$ converge in $L^2$ to the projection of $f$ on the
$\tau$-invariant functions; when the only invariant functions are constants ($\tau$ is \emph{ergodic}), the limit is $\int f\,d\mu$: time averages equal space averages. For the rotation $x\mapsto x+\alpha$ of
Theorem~\ref{thm:kroneckerweyl}, $Te_k=e^{2\pi ik\cdot\alpha}e_k$, and the invariant functions are the constants exactly when $1,\alpha_1,\dots,\alpha_d$ are independent. Birkhoff's pointwise ergodic theorem (1931), proved
in Volume~10, upgrades $L^2$ convergence to almost sure convergence. Von Neumann found his theorem first and told Birkhoff about it, but Birkhoff's paper appeared in December 1931 and von Neumann's only in January 1932, and the question of who should be credited with what is still discussed by historians.

\section{Von Neumann's proof of the Radon--Nikodym theorem}

Volume~6 proved the Radon--Nikodym theorem with the Hahn decomposition. Von Neumann's proof (1940) obtains the density from the Riesz representation theorem, in a few lines.

\begin{theorem}\label{thm:rnvn}
Let $\mu,\nu$ be finite measures on $(E,\mathcal A)$ with $\nu\ll\mu$. There is $h\ge0$ with $\nu(A)=\int_Ah\,d\mu$ for all $A$.
\end{theorem}
\begin{proof}
Let $\rho=\mu+\nu$. The linear form $f\mapsto\int f\,d\nu$ on real $L^2(\rho)$ is bounded: $|\int f\,d\nu|\le\int|f|\,d\rho\le\rho(E)^{1/2}\|f\|_{L^2(\rho)}$. By Theorem~\ref{thm:rieszrep} there is $g\in L^2(\rho)$ with
\[\int f\,d\nu=\int fg\,d\rho\qquad\text{for all }f\in L^2(\rho).\]
With $f=\mathbf 1_A$, $0\le\nu(A)=\int_Ag\,d\rho\le\rho(A)$ for all $A$, so $0\le g\le1$ $\rho$-almost everywhere. Rewrite the identity as $\int f(1-g)\,d\nu=\int fg\,d\mu$. For $B=\{g=1\}$ and $f=\mathbf 1_B$ it gives $\mu(B)=0$,
hence $\nu(B)=0$ by absolute continuity. For $A\in\mathcal A$ and $f=(1+g+\cdots+g^n)\mathbf 1_A$ (bounded, so in $L^2(\rho)$),
\[\int_A(1-g^{n+1})\,d\nu=\int_A(g+g^2+\cdots+g^{n+1})\,d\mu.\]
As $n\to\infty$, the left side tends to $\nu(A\setminus B)=\nu(A)$ by dominated convergence ($g^{n+1}\to0$ off $B$), and the right side increases to $\int_Ah\,d\mu$ with $h=\frac g{1-g}$ (and $h=+\infty$ on the $\mu$-null set $B$, where
we set $h=0$) by monotone convergence.
\end{proof}

The $\sigma$-finite case follows by cutting $E$ into countably many pieces of finite measure. The proof is a model of the Hilbert space method: guess a quantity (here $\int f\,d\nu$) that is a bounded linear form on a well-chosen
$L^2$, and let Riesz produce the object that represents it.

\section{Spaces with a reproducing kernel}

In $L^2$, the value of a function at a point means nothing. Some Hilbert spaces of functions are better: in them, point evaluation is continuous. A \emph{reproducing kernel Hilbert space} (RKHS) on a set $X$ is a Hilbert space $H$ of
functions $X\to\mathbb K$ in which every evaluation $f\mapsto f(x)$ is a bounded linear form. By Theorem~\ref{thm:rieszrep} there is then $k_x\in H$ with $f(x)=\langle f,k_x\rangle$, and the \emph{kernel} of $H$ is
$K(x,y)=k_y(x)=\langle k_y,k_x\rangle$. It is \emph{positive semi-definite}: for points $x_1,\dots,x_n$ and scalars $a_1,\dots,a_n$,
\[\sum_{i,j}a_i\overline{a_j}K(x_j,x_i)=\Bigl\|\sum_ia_ik_{x_i}\Bigr\|^2\ge0.\]

\begin{theorem}[Moore--Aronszajn]\label{thm:aronszajn}
Every positive semi-definite kernel $K$ on a set $X$ (with $K(x,y)=\overline{K(y,x)}$) is the kernel of a unique RKHS of functions on $X$.
\end{theorem}
\begin{proof}
On the span $H_0$ of the functions $K_y=K(\cdot,y)$, define $\langle\sum_ia_iK_{x_i},\sum_jb_jK_{y_j}\rangle=\sum_{i,j}a_i\overline{b_j}K(y_j,x_i)$. It is well defined: with $f=\sum_ia_iK_{x_i}$ and $g=\sum_jb_jK_{y_j}$, it equals
$\sum_j\overline{b_j}f(y_j)$, which depends on $f$ only through its values, and also $\sum_ia_i\overline{g(x_i)}$ (since $K(y_j,x_i)=\overline{K(x_i,y_j)}$), which depends on $g$ only through its values. It is sesquilinear and $\langle f,f\rangle\ge0$ by positivity; and $\langle f,K_x\rangle=f(x)$, the
reproducing property. Cauchy--Schwarz holds for semi-inner products (its proof uses only $\langle f+tg,f+tg\rangle\ge0$), so $|f(x)|=|\langle f,K_x\rangle|\le\langle f,f\rangle^{1/2}K(x,x)^{1/2}$; hence $\langle f,f\rangle=0$ forces $f=0$, and
$H_0$ is an inner product space.

To complete it inside the functions on $X$: if $(f_n)$ is Cauchy in $H_0$, then $(f_n(x))$ is Cauchy for each $x$ by the last inequality; let $f$ be the pointwise limit, and let $H$ be the set of such limits, with
$\|f\|=\lim\|f_n\|$. This is consistent: if $(f_n)$ is Cauchy with $f_n\to0$ pointwise, then $\|f_n\|\to0$. Indeed, with $\|f_n\|\le M$ and $\|f_n-f_m\|\le\varepsilon$ for $n,m\ge N$,
$\|f_n\|^2=\langle f_n,f_n-f_m\rangle+\langle f_n,f_m\rangle$; the first term is at most $M\varepsilon$, and the second, for fixed $n$ with $f_n=\sum_ia_iK_{x_i}$, equals $\sum_ia_i\overline{f_m(x_i)}\to0$ as $m\to\infty$. So
$\|f_n\|^2\le M\varepsilon$ for $n\ge N$. Then $H$, the completion of $H_0$ realised as functions, is a Hilbert space in which $f(x)=\lim f_n(x)=\lim\langle f_n,K_x\rangle=\langle f,K_x\rangle$. Uniqueness: in any RKHS with kernel $K$, the $K_x$
span a dense subspace (a function orthogonal to all $K_x$ vanishes everywhere), on which the inner product is forced.
\end{proof}

The kernel that machine learning uses most is the Gaussian kernel $K(x,y)=e^{-|x-y|^2/2}$ on $\R^d$. It is positive semi-definite because it is a characteristic function: $e^{-|u|^2/2}=\int_{\R^d}e^{i\xi\cdot u}\,d\gamma(\xi)$ with $\gamma$
the standard Gaussian law (Volume~3, Chapter~6), so $\sum_{i,j}a_i\overline{a_j}K(x_j,x_i)=\int|\sum_ia_ie^{-i\xi\cdot x_i}|^2\,d\gamma(\xi)\ge0$. Its RKHS consists of very smooth functions, and the representer theorem
(Exercise~\ref{ex:representer}) reduces learning in this infinite-dimensional space to linear algebra on the $n\times n$ matrix $(K(x_i,x_j))$ of the data. Chapter~17 diagonalises such kernels (Mercer's theorem).

\begin{beauty}
A kernel measures similarity: the number $K(x,y)$ says how much $x$ and $y$ are alike. Moore--Aronszajn turns any reasonable similarity into a geometry, a Hilbert space in which the data points $K_x$ live and where similarity becomes
an inner product. Points that no line can separate in $\R^d$ become separable by a hyperplane there. This is the \emph{kernel trick}, and it is Riesz's theorem at work.
\end{beauty}

\section{Exercises}

\begin{exercise}[$\star\star$]\label{ex:jvn}
\emph{Jordan--von Neumann.} Let $E$ be a real normed space whose norm satisfies the parallelogram law. Show that $\langle x,y\rangle=\frac14(\|x+y\|^2-\|x-y\|^2)$ is an inner product with $\langle x,x\rangle=\|x\|^2$. Deduce that
$L^p$ is a Hilbert space only for $p=2$ (when it has two disjoint sets of positive finite measure).
\end{exercise}

\begin{exercise}[$\star$]\label{ex:legendrefit}
Find the polynomial $p$ of degree at most $2$ minimising $\int_{-1}^1(|x|-p(x))^2\,dx$, and the minimum value.
\end{exercise}

\begin{exercise}[$\star\star$]\label{ex:noclosest}
\emph{No closest point.} On $C([0,1])$ with the sup norm let $\varphi(f)=\int_0^{1/2}f-\int_{1/2}^1f$ and $M=\ker\varphi$. (a) Show that $\|\varphi\|=1$ and that $|\varphi(f)|<\|f\|$ for every $f\ne0$. (b) In any normed space, show that
$d(x,\ker\psi)=\frac{|\psi(x)|}{\|\psi\|}$ for $\psi\in E^*\setminus\{0\}$. (c) Deduce that no $g\notin M$ has a closest point in the closed subspace $M$.
\end{exercise}

\begin{exercise}[$\star\star\star$]\label{ex:vandercorput}
\emph{Van der Corput.} (a) For complex numbers $u_1,\dots,u_N$ (set $u_n=0$ for other $n$) and $1\le H\le N$, prove
\[\Bigl|\sum_nu_n\Bigr|^2\le\frac{N+H-1}{H^2}\sum_{|r|<H}(H-|r|)\gamma(r),\qquad\gamma(r)=\sum_mu_{m+r}\overline{u_m}.\]
(b) Deduce that if, for every integer $r\ge1$, the sequence $(x_{n+r}-x_n)_n$ is equidistributed modulo $1$, then so is $(x_n)$. (c) Show that $n^2\alpha$ is equidistributed modulo $1$ for irrational $\alpha$, and more generally $P(n)$ for
every real polynomial $P$ of degree at least $1$ with irrational leading coefficient (Weyl, 1916).
\end{exercise}

\begin{exercise}[$\star\star$]\label{ex:hardy2}
\emph{The Hardy space.} Let $H^2$ be the space of power series $f(z)=\sum_{n\ge0}a_nz^n$ with $\|f\|^2=\sum_n|a_n|^2<\infty$. (a) Show that each $f\in H^2$ is holomorphic on the unit disc $D$ with $|f(w)|\le\|f\|(1-|w|^2)^{-1/2}$. (b) Show that
$H^2$ is an RKHS on $D$ with kernel $K(z,w)=\frac1{1-z\overline w}$ (Szeg\H o's kernel). (c) Show that $\|f\|^2=\lim_{r\to1}\frac1{2\pi}\int_0^{2\pi}|f(re^{it})|^2\,dt$.
\end{exercise}

\begin{exercise}[$\star\star$]\label{ex:representer}
\emph{The representer theorem.} Let $H$ be a real RKHS on $X$ with kernel $K$, data $(x_i,y_i)\in X\times\R$ for $i\le n$, and $\lambda>0$. Show that $J(f)=\sum_i(f(x_i)-y_i)^2+\lambda\|f\|^2$ has a unique minimiser on $H$, that it lies in
the span of the $K_{x_i}$, and that it is $f=\sum_ia_iK_{x_i}$ with $a=(G+\lambda I)^{-1}y$, where $G=(K(x_i,x_j))_{i,j}$ (kernel ridge regression).
\end{exercise}

\section{Hints}

\begin{hint}{ex:jvn}
Apply the parallelogram law to $x+y$ and $x'$, and to $x-y$ and $x'$, to get $\langle x+x',y\rangle+\langle x-x',y\rangle=2\langle x,y\rangle$. For $L^p$, test $\mathbf 1_A$ and $\mathbf 1_B$.
\end{hint}
\begin{hint}{ex:legendrefit}
$p=c_0P_0+c_1P_1+c_2P_2$ with $c_n=\frac{2n+1}2\int_{-1}^1|x|P_n$; the error is given by Pythagoras.
\end{hint}
\begin{hint}{ex:noclosest}
(a) Approximate $\mathbf 1_{[0,1/2[}-\mathbf 1_{]1/2,1]}$. (b) $\psi(x)=\psi(x-m)$ for $m\in\ker\psi$. (c) A closest point would make $\varphi$ attain its norm.
\end{hint}
\begin{hint}{ex:vandercorput}
(a) Write $H\sum_nu_n=\sum_n\sum_{h<H}u_{n+h}$, with $N+H-1$ non-zero values of $n$, and apply Cauchy--Schwarz. (b) Weyl's criterion, with $u_n=e^{2\pi ikx_n}$. (c) $(n+r)^2\alpha-n^2\alpha=2r\alpha n+r^2\alpha$; induct on the degree.
\end{hint}
\begin{hint}{ex:hardy2}
$f(w)=\langle f,k_w\rangle$ with $k_w(z)=\sum_n\overline w^nz^n$. For (c), Parseval on the circle of radius $r$.
\end{hint}
\begin{hint}{ex:representer}
Split $f=v+w$ with $v$ in the span $V$ of the $K_{x_i}$ and $w\perp V$; then $f(x_i)=v(x_i)$.
\end{hint}

\section{Solutions}

\begin{solution}{ex:jvn}
\idea The candidate is forced by the polarisation identity. What must be shown is additivity in $x$, and the parallelogram law is a relation between four norms, which is exactly the size of an additivity statement.
\writeup Symmetry and $\langle x,x\rangle=\|x\|^2$ are clear. The parallelogram law applied to $x+y$ and $x'$, then to $x-y$ and $x'$, gives
$\|x+x'+y\|^2+\|x-x'+y\|^2=2\|x+y\|^2+2\|x'\|^2$ and $\|x+x'-y\|^2+\|x-x'-y\|^2=2\|x-y\|^2+2\|x'\|^2$. Subtracting and dividing by $4$:
$\langle x+x',y\rangle+\langle x-x',y\rangle=2\langle x,y\rangle$. With $x'=x$ (and $\langle0,y\rangle=0$), $\langle2x,y\rangle=2\langle x,y\rangle$. For any $a,b$, put $x=\frac{a+b}2$, $x'=\frac{a-b}2$: $\langle a,y\rangle+\langle b,y\rangle=2\langle\frac{a+b}2,y\rangle=\langle a+b,y\rangle$.
So $\langle\cdot,y\rangle$ is additive, hence $\langle qx,y\rangle=q\langle x,y\rangle$ for rational $q$, and for real $q$ by continuity of the norm. Positivity is $\langle x,x\rangle=\|x\|^2$. For $L^p$, if $A,B$ are disjoint with
$\mu(A)=\mu(B)=1$, then $\|\mathbf 1_A\pm\mathbf 1_B\|_p^2=2^{2/p}$ and the law reads $2\cdot2^{2/p}=4$, so $p=2$; if $\mu(A)=a$, $\mu(B)=b$, rescale the indicators.
\end{solution}

\begin{solution}{ex:legendrefit}
\idea Best approximation in $L^2$ is orthogonal projection, and with an orthogonal basis the projection is read off coefficient by coefficient. Symmetry kills the odd Legendre polynomials.
\writeup By Theorem~\ref{thm:projection} and Proposition~\ref{prop:legendre}, $p=\sum_{n\le2}c_nP_n$ with $c_n=\frac{\int|x|P_n}{\int P_n^2}=\frac{2n+1}2\int_{-1}^1|x|P_n$. $c_1=0$ by parity, $c_0=\frac12$, and
$c_2=\frac52\cdot2\int_0^1x\frac{3x^2-1}2\,dx=\frac52(\frac34-\frac12)=\frac58$. So $p(x)=\frac12+\frac58\cdot\frac{3x^2-1}2=\frac3{16}+\frac{15}{16}x^2$. By Pythagoras the minimum is
$\int x^2-c_0^2\cdot2-c_2^2\cdot\frac25=\frac23-\frac12-\frac5{32}=\frac1{96}$.
\end{solution}

\begin{solution}{ex:noclosest}
\idea In a Hilbert space the closest point exists because of the parallelogram law. In $C([0,1])$ the closest point would have to make a linear form attain its norm, and this form only reaches its norm on a discontinuous
function.
\writeup (a) $|\varphi(f)|\le\int_0^1|f|\le\|f\|$. Continuous functions $f_n$, equal to $1$ on $[0,\frac12-\frac1n]$, to $-1$ on $[\frac12+\frac1n,1]$ and linear in between, have $\|f_n\|=1$ and $\varphi(f_n)=1-\frac1n$, so $\|\varphi\|=1$. If
$\|f\|=1$ and $\varphi(f)=1$, then $\int_0^{1/2}(1-f)+\int_{1/2}^1(1+f)=0$ with continuous non-negative integrands, so $f=1$ on $[0,\frac12]$ and $f=-1$ on $[\frac12,1]$, impossible at $\frac12$. Replacing $f$ by $-f$ handles
$\varphi(f)=-1$; and scaling gives $|\varphi(f)|<\|f\|$ for all $f\ne0$. (b) For $m\in\ker\psi$, $|\psi(x)|=|\psi(x-m)|\le\|\psi\|\|x-m\|$. Conversely, given $\varepsilon$, pick $u$ with $\|u\|=1$ and $|\psi(u)|>\|\psi\|-\varepsilon$;
$m=x-\frac{\psi(x)}{\psi(u)}u$ is in $\ker\psi$ and $\|x-m\|=\frac{|\psi(x)|}{|\psi(u)|}$. (c) If $m\in M$ were closest to $g$, then $\|g-m\|=|\varphi(g)|=|\varphi(g-m)|$, contradicting (a) since $g-m\ne0$.
\end{solution}

\begin{solution}{ex:vandercorput}
\idea To show that a sum of unimodular numbers is small, square it: the square involves the correlations $u_{n+r}\overline{u_n}$, which may be easier to understand. For $u_n=e^{2\pi ikx_n}$ they are exponentials of the
differences $x_{n+r}-x_n$, and differencing lowers the degree of a polynomial.
\writeup (a) For each $m$, $u_m$ appears in $\sum_{h<H}u_{n+h}$ for exactly $H$ values of $n$, so $H\sum_nu_n=\sum_n\sum_{0\le h<H}u_{n+h}$, where the inner sum vanishes unless $1-H<n\le N$: at most $N+H-1$ values. By
Cauchy--Schwarz,
\[H^2\Bigl|\sum_nu_n\Bigr|^2\le(N+H-1)\sum_n\Bigl|\sum_{h<H}u_{n+h}\Bigr|^2=(N+H-1)\sum_{h,h'<H}\sum_nu_{n+h}\overline{u_{n+h'}},\]
and $\sum_nu_{n+h}\overline{u_{n+h'}}=\gamma(h-h')$, while each $r$ with $|r|<H$ equals $h-h'$ for exactly $H-|r|$ pairs. (b) Fix $k\ne0$ and $u_n=e^{2\pi ikx_n}$ for $n\le N$. Then $\gamma(0)=N$, $\gamma(-r)=\overline{\gamma(r)}$, and
$\gamma(r)=\sum_{m=1}^{N-r}e^{2\pi ik(x_{m+r}-x_m)}$ for $r\ge1$, so by (a)
\[\Bigl|\frac1N\sum_{n\le N}u_n\Bigr|^2\le\frac{N+H-1}{NH}+\frac{2(N+H-1)}{NH}\sum_{r=1}^{H-1}\frac{|\gamma(r)|}N.\]
For fixed $H$ and $N\to\infty$, each $\frac{\gamma(r)}N\to0$ by Weyl's criterion (Volume~4) for the sequence $x_{m+r}-x_m$, so the $\limsup$ is at most $\frac1H$. As $H$ is arbitrary, $\frac1N\sum_ne^{2\pi ikx_n}\to0$ for all $k\ne0$, and
$(x_n)$ is equidistributed by Weyl's criterion. (c) Induct on the degree $d$. For $d=1$, $P(n)=\alpha n+\beta$, and $\frac1N\sum_ne^{2\pi ik(\alpha n+\beta)}=e^{2\pi ik\beta}\frac1N\sum_ne^{2\pi ik\alpha n}\to0$ for $k\ne0$ (Volume~4): equidistributed. If $d\ge2$ and the leading
coefficient $\alpha_d$ is irrational, then for each $r\ge1$, $P(n+r)-P(n)$ is a polynomial in $n$ of degree $d-1$ with leading coefficient $rd\alpha_d$, irrational; by induction it is equidistributed, and (b) applies. For
$n^2\alpha$ the differences are $2r\alpha n+r^2\alpha$.
\end{solution}

\begin{solution}{ex:hardy2}
\idea $H^2$ is $\ell^2$ in disguise, with the coefficients as coordinates. Evaluation at $w$ is the inner product with the sequence $(\overline w^n)$, which is square summable exactly when $|w|<1$.
\writeup (a) For $|w|<1$, Cauchy--Schwarz gives $\sum|a_n||w|^n\le\|f\|(\sum|w|^{2n})^{1/2}=\|f\|(1-|w|^2)^{-1/2}$: the series converges absolutely, normally on each disc $|w|\le r<1$, so $f$ is holomorphic on $D$ (Volume~4, Chapter~9) and
satisfies the bound. (b) $H^2$ is isometric to $\ell^2$, hence a Hilbert space, and $f(w)=\sum a_nw^n=\langle f,k_w\rangle$ with $k_w=\sum_n\overline w^nz^n\in H^2$. So evaluations are bounded, and $K(z,w)=k_w(z)=\sum_n(z\overline w)^n=\frac1{1-z\overline w}$.
(c) The function $t\mapsto f(re^{it})=\sum_na_nr^ne^{int}$ has Fourier coefficients $a_nr^n$ (the series converges uniformly in $t$), so by Parseval $\frac1{2\pi}\int|f(re^{it})|^2\,dt=\sum_n|a_n|^2r^{2n}$, which increases to $\|f\|^2$ by
monotone convergence as $r\uparrow1$.
\end{solution}

\begin{solution}{ex:representer}
\idea The data only see $f$ through its values at the $x_i$, which are inner products with the $K_{x_i}$. The component of $f$ orthogonal to them is invisible to the data and costs norm, so it is zero at the optimum.
The problem becomes finite-dimensional.
\writeup Write $f=v+w$ with $v\in V=\operatorname{span}(K_{x_i})$ and $w\perp V$ (Corollary~\ref{cor:orthdecomp}; $V$ is finite-dimensional, hence closed). Then $f(x_i)=\langle f,K_{x_i}\rangle=v(x_i)$ and $\|f\|^2=\|v\|^2+\|w\|^2$, so
$J(f)=J(v)+\lambda\|w\|^2$: a minimiser lies in $V$. On $V$, $J$ is continuous and $J(f)\ge\lambda\|f\|^2$, so it attains its minimum on a large ball, a compact set in finite dimension. It is unique because $J$ is strictly
convex ($f\mapsto\|f\|^2$ is strictly convex by the parallelogram law, and the other terms are convex). For $f=\sum_ia_iK_{x_i}$, $f(x_j)=(Ga)_j$ and $\|f\|^2=a^TGa$, so $J=|Ga-y|^2+\lambda a^TGa$, whose gradient in $a$ is
$2G(Ga-y)+2\lambda Ga=2G((G+\lambda I)a-y)$. It vanishes at $a=(G+\lambda I)^{-1}y$ ($G$ is positive semi-definite, so $G+\lambda I$ is invertible). A convex differentiable function is minimal where its gradient vanishes, so
this $a$ gives the minimiser.
\end{solution}

\section*{Sources and further reading}

\begin{itemize}
\item E.~Schmidt, ``\"Uber die Aufl\"osung linearer Gleichungen mit unendlich vielen Unbekannten'', \emph{Rend. Circ. Mat. Palermo} 25 (1908), 53--77. The geometry of $\ell^2$.
\item F.~Riesz, ``Sur les syst\`emes orthogonaux de fonctions'', and E.~Fischer, ``Sur la convergence en moyenne'', \emph{C. R. Acad. Sci. Paris} 144 (1907); F.~Riesz and M.~Fr\'echet on the representation of linear forms, same volume.
\item J.~von Neumann, ``Proof of the quasi-ergodic hypothesis'', \emph{Proc. Nat. Acad. Sci.} 18 (1932), 70--82, and G.~D.~Birkhoff, ``Proof of the ergodic theorem'', \emph{Proc. Nat. Acad. Sci.} 17 (1931), 656--660; the priority question is examined in J.~D.~Zund, ``George David Birkhoff and John von Neumann: a question of priority and the ergodic theorems, 1931--1932'', \emph{Historia Math.} 29 (2002), 138--153; and ``On rings of operators III'', \emph{Ann. of Math.} 41 (1940), 94--161, for the proof of Radon--Nikodym.
\item N.~Aronszajn, ``Theory of reproducing kernels'', \emph{Trans. Amer. Math. Soc.} 68 (1950), 337--404.
\item H.~Weyl, ``\"Uber die Gleichverteilung von Zahlen mod.~Eins'', \emph{Math. Ann.} 77 (1916), 313--352; L.~Kuipers and H.~Niederreiter, \emph{Uniform Distribution of Sequences}, Wiley, 1974 (van der Corput's inequality).
\item N.~Young, \emph{An Introduction to Hilbert Space}, Cambridge, 1988; H.~Brezis, \emph{Functional Analysis}, Springer, 2011, Chapter~5.
\item B.~Sch\"olkopf and A.~Smola, \emph{Learning with Kernels}, MIT Press, 2002. Kernels in machine learning.
\item C.~Reid, \emph{Hilbert}, Springer, 1970. The life of Hilbert and of G\"ottingen mathematics.
\item For the stories: Fredholm's insurance work is in his biography in the MacTutor History of Mathematics archive (St Andrews); Holmgren's report
to Hilbert is in Reid's book above; J.~von Neumann, ``Allgemeine Eigenwerttheorie Hermitescher Funktionaloperatoren'', \emph{Math.\ Ann.} 102 (1929), 49--131, axiomatises Hilbert space;
L.~A.~Steen, ``Highlights in the history of spectral theory'', \emph{Amer.\ Math.\ Monthly} 80 (1973), 359--381.
\end{itemize}
